\subsection{Runtime ablations}
\label{sec:results-practical}

\Cref{tab:runtime-ablations} changes one setting at a time and reports fitting
time relative to the selected configuration of the same estimator. Each ratio
is the median over datasets of the per-dataset ratio, given as the range over
the four task groups (classification and regression, real and synthetic).
Ratios above 1 are slower. \Cref{app:complexity-details} links each row to the
work terms it affects.

\begin{table}[!htbp]
  \centering
  \footnotesize
  \setlength{\tabcolsep}{4pt}
  \begin{tabular}{@{}lccccc@{}}
    \toprule
     & \multicolumn{2}{c}{CIT} & \multicolumn{3}{c}{CIF} \\
    \cmidrule(lr){2-3}\cmidrule(l){4-6}
    \tablehead{Change} & Time ratio & \shortstack{Top-10 F1\\change} & Time ratio & \shortstack{Downstream\\change} & \shortstack{Recovery\\change} \\
    \midrule
    Disable adaptive stopping & 0.88 to 0.99 & 0 & 0.96 to 1.00 & 0 & 0 \\
    Auto budget, adaptive stopping & 1.18 to 1.56 & $\le 0.016$ & 1.30 to 2.0 & & $\le 0.030$ \\
    Auto budget, no stopping & 3.8 to 8.2 & $\le 0.016$ & 5.8 to 15 & & $\le 0.025$ \\
    Disable threshold scanning & 2.7 to 34 & $\le 0.033$ & 7.9 to 30 & $\le 0.003$ & $\le 0.010$ \\
    Disable feature scanning & 0.81 to 3.0 & $\le 0.058$ & 1.0 to 5.2 & $\le 0.003$ & $\le 0.008$ \\
    Exact threshold search & 0.91 to 1.10 & $\le 0.004$ & 1.3 to 19 & $\le 0.009$ & $\le 0.018$ \\
    Histogram with 32 bins & & & 0.05 to 0.14 & $\le 0.008$ & $\le 0.013$ \\
    Disable feature muting & 0.97 to 0.98 & 0 & 0.91 to 1.00 & $\le 0.001$ & $\le 0.003$ \\
    Disable bootstrap sampling & & & 0.94 to 1.69 & $\le 0.012$ & $\le 0.011$ \\
    Max-type Stage~B rule & & & 0.14 to 0.61 & $\le 0.005$ & $\le 0.007$ \\
    Remove both corrections & 0.52 to 1.8 & $\le 0.063$ & 0.03 to 0.34 & $\le 0.004$ & $\le 0.011$ \\
    \bottomrule
  \end{tabular}
  \caption{Runtime ablations of CIT and CIF on the runtime panel of 23 datasets
  (15 classification and 8 regression; 9 real and 14 synthetic), five seeds per
  dataset. Time ratios are medians over datasets of the per-dataset ratio to
  the selected configuration, given as the range over the four task groups.
  Top-10 F1 is the harmonic mean of precision and recall of the ten
  highest-ranked features against the planted informative set. Score columns
  give the largest absolute change across task groups in mean top-10 F1 for
  CIT, and in the dataset-mean downstream score and synthetic feature recovery
  for CIF; bounds are rounded upward to 0.001. The auto-budget runs record
  recovery but no downstream score. Empty cells are settings that do not apply
  to, or were not ablated for, that estimator. Every time is the second of two
  identical fits, so compilation is excluded.}
  \label{tab:runtime-ablations}
\end{table}

At the benchmark's minimum budget, adaptive stopping is inert by construction
(\cref{sec:theory-adaptive}). Switching it off leaves every tree unchanged, and
the time ratios near 1 measure only the cost of the check. The result does not
depend on how forest trees are scheduled, since a forest with one worker and
the full thread pool inside each tree gives 0.95 to 1.01.

At the auto budget, where the rule can act, adaptive stopping brings the tree
to 0.14 to 0.37 and the forest to 0.11 to 0.25 of their no-stopping time, with
recovery changes of at most 0.030.

Of the runtime controls, the threshold scan has the largest effect. Testing
candidate thresholds in order of weighted child impurity and stopping at the
first rejection makes trees 2.7 to 34 times and forests 7.9 to 30 times faster
than testing every candidate. The candidate count has the next largest effect.
For the forest, exact threshold search is 1.3 to 19 times slower than the
256-bin histogram, and a 32-bin histogram is 7 to 21 times faster, with
downstream changes of at most 0.009. Feature scanning matters most on the real regression datasets. The
max-type rule fits the forest in 0.14 to 0.61 of the reference time.

Removing the two corrections changes the statistical rule as well as the
fitting time. It deepens trees from a mean of 6.5 and 7.2 levels to 10.7 and
13.4, and moves top-10 recovery by up to 0.063. The selected configuration
keeps both corrections.

The pattern holds on the 14 Benchmark-size datasets, with up to 20,000
observations and 5,000 features. Switching adaptive stopping off changes tree
time by 0.80 to 1.01 (medians 0.97 in classification and 0.88 in regression)
and forest time by 0.99 to 1.34 (medians 1.00 and 1.03). Threshold scanning
saves a median 17 times in classification and 45 times on spam, and removing
the corrections runs in 0.01 to 1.26 of the reference time. The slowest single
tree is gisette at 29 minutes, where the feature correction over 5,000 features
sets a large minimum budget for each feature test.
